% Źródła: Eves s. 295--298 (§7.3).

\section{Trójkąty graniczne i absolutna długość} % lang-pl
\section{Triangoli limite e lunghezza assoluta} % lang-it
\label{sekcja_trojkaty_graniczne}

Łobaczewski, Bolyai i Gauss odkryją, że na płaszczyźnie hiperbolicznej istnieje naturalna jednostka długości, i zrobią to za pomocą figury, której nie ma w geometrii euklidesowej. % lang-pl
W tej sekcji poznamy trójkąty graniczne, kąt równoległości jako funkcję odległości i zdumiewający wniosek stąd płynący. % lang-pl
Lobachevskij, Bolyai e Gauss scopriranno che nel piano iperbolico esiste un'unità naturale di lunghezza, e lo faranno con una figura che nella geometria euclidea non esiste. % lang-it
In questa sezione conosceremo i triangoli limite, l'angolo di parallelismo come funzione della distanza e la sorprendente conseguenza che ne deriva. % lang-it

\begin{definition}
[trójkąt graniczny] % lang-pl
[triangolo limite] % lang-it
	Dwie półproste równoległe wraz z odcinkiem łączącym ich początki. % lang-pl
	Odcinek nazywamy bokiem skończonym, a kąty przy jego końcach -- kątami trójkąta granicznego. % lang-pl
	Due semirette parallele insieme al segmento che ne congiunge le origini. % lang-it
	Il segmento si chiama lato finito, mentre gli angoli ai suoi estremi sono gli angoli del triangolo limite. % lang-it
\index{trójkąt!graniczny} % lang-pl
\index{triangolo!limite} % lang-it
\end{definition}

Eves \cite[s. 295]{eves1_1972}. % lang-pl
Eves \cite[p. 295]{eves1_1972}. % lang-it

\begin{proposition}
	Kąt zewnętrzny trójkąta granicznego jest większy od kąta wewnętrznego do niego nieprzyległego. % lang-pl
	Jeśli boki skończone dwóch trójkątów granicznych są równe i jeden kąt jednego jest równy kątowi drugiego, to pozostałe kąty też są równe; jeśli dwa kąty jednego są równe dwóm kątom drugiego, to boki skończone są równe. % lang-pl
	L'angolo esterno di un triangolo limite è maggiore dell'angolo interno non adiacente. % lang-it
	Se i lati finiti di due triangoli limite sono uguali e un angolo dell'uno è uguale a un angolo dell'altro, allora anche gli altri angoli sono uguali; se due angoli dell'uno sono uguali a due angoli dell'altro, allora i lati finiti sono uguali. % lang-it
\end{proposition}

Dowody: Eves \cite[s. 295--296]{eves1_1972}. % lang-pl
Dimostrazioni: Eves \cite[p. 295--296]{eves1_1972}. % lang-it

Wynika stąd, że kąt równoległości w punkcie $P$ względem prostej $AB$ zależy tylko od odległości $h$ punktu $P$ od $AB$ i maleje, gdy $h$ rośnie \cite[s. 296--297]{eves1_1972}. % lang-pl
Łobaczewski oznaczy go przez $\Pi(h)$. % lang-pl
Eves zapowiada, że przy odpowiednim wyborze jednostki długości zachodzi $\Pi(h) = 2\arctan e^{-h}$ \cite[s. 297]{eves1_1972}; $\Pi(h)$ maleje od $\tfrac{\pi}{2}$ do $0$, gdy $h$ rośnie od $0$ do $\infty$ (Eves \cite[s. 298, zad. 9]{eves1_1972}). % lang-pl
Ne segue che l'angolo di parallelismo nel punto $P$ rispetto alla retta $AB$ dipende soltanto dalla distanza $h$ di $P$ da $AB$ e diminuisce al crescere di $h$ \cite[p. 296--297]{eves1_1972}. % lang-it
Lobachevskij lo indicherà con $\Pi(h)$. % lang-it
Eves annuncia che, con un'opportuna scelta dell'unità di lunghezza, vale $\Pi(h) = 2\arctan e^{-h}$ \cite[p. 297]{eves1_1972}; $\Pi(h)$ decresce da $\tfrac{\pi}{2}$ a $0$ quando $h$ cresce da $0$ a $\infty$ (Eves \cite[p. 298, es. 9]{eves1_1972}). % lang-it

Konsekwencja jest zaskakująca. % lang-pl
W geometrii euklidesowej kąty są \emph{absolutne} (istnieje naturalna jednostka: kąt prosty), ale długości są \emph{względne} -- gdyby zgubić wzorzec metra, nie da się go odtworzyć z samej geometrii. % lang-pl
W geometrii Łobaczewskiego również długości są absolutne: każdemu kątowi równoległości odpowiada określona odległość $h$, więc z jednostki kąta da się odtworzyć jednostkę długości \cite[s. 297]{eves1_1972}. % lang-pl
La conseguenza è sorprendente. % lang-it
In geometria euclidea gli angoli sono \emph{assoluti} (esiste un'unità naturale: l'angolo retto), ma le lunghezze sono \emph{relative} -- se si perdesse il campione del metro, non si potrebbe ricostruirlo dalla sola geometria. % lang-it
In geometria di Lobachevskij anche le lunghezze sono assolute: a ogni angolo di parallelismo corrisponde una distanza $h$ ben determinata, quindi da un'unità angolare si ricava un'unità di lunghezza \cite[p. 297]{eves1_1972}. % lang-it

Eves dodaje, że „w małej skali'' geometria Łobaczewskiego przybliża euklidesową (zad. 9c, s. 298); pytanie, dlaczego Biuro Miar przechowuje wzorzec długości, a nie kąta, jest jego zadaniem 10 \cite[s. 298]{eves1_1972}. % lang-pl
Eves aggiunge che «in piccolo» la geometria di Lobachevskij approssima quella euclidea (es. 9c, p. 298); la domanda sul perché l'Ufficio dei Pesi e delle Misure conservi un campione di lunghezza e non di angolo è il suo esercizio 10 \cite[p. 298]{eves1_1972}. % lang-it
